\subsection{Consistency Conditions for the SBC/F3 Infrared Kernel}

For the SBC/F3 sector to remain observationally viable, the memory kernel must satisfy a set of physical and phenomenological consistency conditions.

First, causality requires that the effective response at redshift \(z\) only depends on the accessible causal history:
\begin{equation}
K(z,z')=0
\qquad
{\rm for}
\qquad
z'<z .
\end{equation}

Second, the kernel must be infrared dominated. This requires exponential suppression away from the late-time sector,
\begin{equation}
K(z,z') \propto \exp[-k_0 |z'-z|],
\qquad
k_0>0 .
\end{equation}

Third, the amplitude of the response must remain perturbative:
\begin{equation}
|\epsilon| \ll 1 .
\end{equation}

The observationally preferred value,
\begin{equation}
\epsilon\simeq 0.06,
\end{equation}
satisfies this condition and explains why the high-\(\ell\) CMB acoustic structure, BAO distance relations, and Pantheon+ luminosity distances remain stable.

Fourth, the homogeneous background deformation must remain strongly suppressed. The joint BAO and Pantheon+ analysis constrains the background coupling to
\begin{equation}
\alpha_{\rm bg}\simeq 0,
\end{equation}
which implies that the viable SBC/F3 sector cannot be a large background-level modification of \(\Lambda\)CDM.

The allowed structure is therefore
\begin{equation}
H_{\rm SBC}(z)\simeq H_{\Lambda{\rm CDM}}(z),
\end{equation}
while perturbative observables may receive infrared corrections:
\begin{equation}
X_{\rm SBC}(z,k)=X_{\Lambda{\rm CDM}}(z,k)
\left[1+F_{\rm SBC}(z,k)\right].
\end{equation}

This leads to the observationally preferred hierarchy
\begin{equation}
|F_{\rm background}|
\ll
|F_{\rm perturbative}|
\ll
1 .
\end{equation}

These conditions explain why the SBC/F3 sector can remain compatible with background probes while still allowing controlled perturbative infrared deviations.
